\section{Impulse limit and Doppler width}\label{app:D.32}

A nucleus of mass $M_\kappa$ at position $\vect{X}$ with momentum $\vect{p}$ takes up the momentum $\vect{q}$ through $\op{O} = \ee^{\ii\vect{q}\cdot\vect{X}}$. We derive the moments of the self response of this nucleus for any potential, the impulse limit of the self response and the width of the self response in a harmonic crystal. Throughout, $E_R = q^2/2M_\kappa$ and $p_{\hat q} = \hat{\vect{q}}\cdot\vect{p}$.

\subsection{The self correlation}

The self part of \cref{eq:11.3-exact} for one nucleus is $c_\kappa^2\Omega^{-1}\int\dd t\,\ee^{\ii\omega t}I(t)$ with $I(t) = \langle\ee^{-\ii\vect{q}\cdot\vect{X}(t)}\ee^{\ii\vect{q}\cdot\vect{X}}\rangle$. Write $\ee^{-\ii\vect{q}\cdot\vect{X}(t)} = \ee^{\ii Ht}\ee^{-\ii\vect{q}\cdot\vect{X}}\ee^{-\ii Ht}$ and move the kick through the second evolution operator,
\begin{equation*}
  \ee^{-\ii\vect{q}\cdot\vect{X}}\,\ee^{-\ii Ht}\,\ee^{\ii\vect{q}\cdot\vect{X}} = \ee^{-\ii H_{\vect{q}}t} , \qquad H_{\vect{q}} = \ee^{-\ii\vect{q}\cdot\vect{X}}H\,\ee^{\ii\vect{q}\cdot\vect{X}} .
\end{equation*}
The translation $\ee^{-\ii\vect{q}\cdot\vect{X}}\vect{p}\,\ee^{\ii\vect{q}\cdot\vect{X}} = \vect{p} + \vect{q}$ changes the kinetic energy $\vect{p}^2/2M_\kappa$ and nothing else, provided the potential energy depends on $\vect{X}$ but not on $\vect{p}$. Every Coulomb interaction depends on positions only, not on momenta. Then
\begin{equation}
  I(t) = \bigl\langle\ee^{\ii Ht}\,\ee^{-\ii(H + D)t}\bigr\rangle , \qquad D = H_{\vect{q}} - H = E_R + \frac{\vect{q}\cdot\vect{p}}{M_\kappa} ,
  \label{eq:D.32-I}
\end{equation}
exactly, for any potential and any stationary state. The operator $D$ is the energy that a kick $\vect{q}$ would give a free particle of momentum $\vect{p}$.

\subsection{Moments of the recoil spectrum}

The moments of the spectrum are the derivatives of $I$ at $t = 0$, $\int\dd\omega\,\omega^kS/\int\dd\omega\,S = \ii^k\,\dd^kI/\dd t^k|_0$ with the sign of \cref{eq:4.2-Sab}. The first derivative gives $\langle(H + D) - H\rangle = \langle D\rangle = E_R$, since $\langle\vect{p}\rangle = 0$; the first moment $E_R$ is \cref{eq:5.6-sum}. The second derivative gives $\langle H^2 - 2H(H + D) + (H + D)^2\rangle = \langle D^2\rangle + \langle[D, H]\rangle$, and the average of a commutator with $H$ vanishes in a stationary state. Hence the variance
\begin{equation}
  \sigma_\omega^2 = \langle D^2\rangle - E_R^2 = \frac{q^2\langle p_{\hat q}^2\rangle}{M_\kappa^2} = 4E_R\,E^{\mathrm{kin}}_\kappa(\hat{\vect{q}}) , \qquad E^{\mathrm{kin}}_\kappa(\hat{\vect{q}}) = \frac{\langle p_{\hat q}^2\rangle}{2M_\kappa} ,
  \label{eq:D.32-var}
\end{equation}
holds at any temperature and for any potential local in $\vect{X}$. The width of the response measures the kinetic energy of the nucleus along $\hat{\vect{q}}$. At $T = 0$, with $|0\rangle$ the ground state, the third moment is $\langle0|(H - E_0 + D)^3|0\rangle$. Every term with $H - E_0$ at either end vanishes, which leaves $\langle D^3\rangle + \langle D(H - E_0)D\rangle$. The term $\langle D^3\rangle$ is $E_R^3 + 3E_R\sigma_\omega^2$, because $\langle p_{\hat q}^3\rangle = 0$ on a site with inversion symmetry or by time reversal. The term $\langle D(H - E_0)D\rangle$ is half the double commutator $\langle[D,[H, D]]\rangle$, and with $[V, p_{\hat q}] = \ii\,\partial_{\hat q}V$ the double commutator gives $(q^2/2M_\kappa^2)\langle\partial_{\hat q}^2V\rangle$, the mean curvature of the potential along $\hat{\vect{q}}$. Subtracting the powers of the lower moments gives the third cumulant
\begin{equation}
  c_3 = \frac{q^2}{2M_\kappa^2}\,\bigl\langle\partial_{\hat q}^2V\bigr\rangle .
  \label{eq:D.32-c3}
\end{equation}
For a nucleus whose electrons follow it, at $n_B = 0$, $V$ is the Born--Oppenheimer surface of \cref{sec:5.1}. With the full Coulomb Hamiltonian instead, $V$ contains the attraction of the electrons, whose curvature at the nucleus is the electron density there. The third moment then includes the electronic excitations that the kick causes, and \cref{sec:12.5} treats it.

\subsection{The moments at any temperature}

In thermal equilibrium the response is the average over the initial eigenstates $|i\rangle$ of $H$ with the Boltzmann weights $p_i$. From one eigenstate, the moments of the spectrum of $\omega = E_f - E_i$ are $\sum_f|\langle f|\ee^{\ii\vect{q}\cdot\vect{X}}|i\rangle|^2(E_f - E_i)^k = \langle i|(A_i + D)^k|i\rangle$, with $A_i = H - E_i$, by the same translation of $H$ as above. For a Hamiltonian invariant under time reversal the eigenfunctions can be chosen real, and then $\langle i|p_{\hat q}|i\rangle = \langle i|p_{\hat q}^3|i\rangle = 0$: in a real state the expectation value of $p_{\hat q} = -\ii\partial_{\hat q}$, or of its cube, is $\ii$ times a real integral, while an expectation value of a Hermitian operator must be real. The mean is therefore $E_R$ from every $|i\rangle$, and the central moments are $\langle i|(A_i + \delta D)^k|i\rangle$ with $\delta D = (q/M_\kappa)p_{\hat q}$. Since $A_i|i\rangle = 0$, every term with $A_i$ at either end vanishes, exactly as for the ground state. The second moment is $\langle i|\delta D^2|i\rangle$. The third is $\langle i|\delta D^3|i\rangle + \langle i|\delta D\,A_i\,\delta D|i\rangle$; the first term vanishes, and the second is $\frac12\langle i|[\delta D,[H,\delta D]]|i\rangle = (q^2/2M_\kappa^2)\langle i|\partial_{\hat q}^2V|i\rangle$. The fourth follows by the steps of \cref{sec:11.5} with $|i\rangle$ in place of the ground state and $A_i$ in place of $A$. Averaging over $i$ with the weights $p_i$ gives the moments \cref{eq:11.5-moments} with thermal averages in place of ground-state averages: all five moments hold at any temperature. For one harmonic oscillator of frequency $\omega_0$ the fourth cumulant that results is $E_R\omega_0^3\coth(\omega_0/2T)$, the even case of \cref{eq:D.32-ck} below.

\subsection{The impulse approximation}

If $D$ commuted with $H$, \cref{eq:D.32-I} would be $I(t) = \langle\ee^{-\ii Dt}\rangle$. For a free particle $I(t) = \langle\ee^{-\ii Dt}\rangle$ is exact. For a bound nucleus the commutator $[H, D] = \ii(q/M_\kappa)\,\partial_{\hat q}V$ is the force along $\hat{\vect{q}}$, and dropping it is the impulse approximation: the nucleus is treated as free during the kick. Written as an average over the momentum distribution $\mathfrak{f}_\kappa(\vect{p})$ of the nucleus, normalized to one, the impulse approximation is
\begin{equation}
  I_{\mathrm{IA}}(t) = \int\dd^3p\;\mathfrak{f}_\kappa(\vect{p})\,\ee^{-\ii(E_R + \vect{q}\cdot\vect{p}/M_\kappa)t} , \qquad
  \int\dd t\;\ee^{\ii\omega t}I_{\mathrm{IA}}(t) = 2\pi\int\dd^3p\;\mathfrak{f}_\kappa(\vect{p})\,\delta\Bigl(\omega - E_R - \frac{\vect{q}\cdot\vect{p}}{M_\kappa}\Bigr) .
  \label{eq:D.32-IA}
\end{equation}
The mean and variance of the impulse spectrum are those of \cref{eq:D.32-var}, exactly, and its third cumulant vanishes. The approximation errs first in $c_3$. In the scaled variable $(\omega - E_R)/\sigma_\omega$ the omitted skewness is $c_3/\sigma_\omega^3$. With $c_3 \propto q^2$ and $\sigma_\omega \propto q$, the omitted skewness falls as $1/q$, while every cumulant that \cref{eq:D.32-IA} keeps scales with the width. The impulse approximation is the leading term of an expansion in $1/q$ at fixed $(\omega - E_R)/\sigma_\omega$.

\subsection{The harmonic crystal}

In a harmonic crystal $\mathfrak{f}_\kappa$ is Gaussian, and its width follows from \cref{eq:5.6-C}: $\langle p^\alpha p^\beta\rangle = M_\kappa^2\langle\dot u^\alpha\dot u^\beta\rangle$ is the mode sum of $(M_\kappa/2N)\,\omega_{\vect{Q}\nu}e^\alpha_{\kappa\nu}e^{\beta*}_{\kappa\nu}\coth(\omega_{\vect{Q}\nu}/2T)$, which on a cubic site gives the kinetic energy $E^{\mathrm{kin}}_\kappa$ of \cref{sec:11.4}, $\langle\omega\rangle/4$ at $T = 0$. The impulse spectrum \cref{eq:D.32-IA} is then a Gaussian of mean $E_R$ and variance $E_R\langle\omega\coth(\omega/2T)\rangle$. The same Gaussian follows from the generating formula \cref{eq:5.6-gen} directly. The exponent of the generating formula, $\vect{q}\cdot\mathcal{C}_{\kappa\kappa}(t)\cdot\vect{q}$, expanded to second order in $t$, is $2W_\kappa - \ii E_Rt - \frac12E_R\langle\omega\coth(\omega/2T)\rangle t^2$; the constant cancels the Debye--Waller factor, and the time integral is Gaussian. This Gaussian is the saddle-point evaluation of Campbell-Deem and collaborators~\cite[Eqs.~(37)--(39) and App.~B]{campbelldeem2022dark}. The $k$-th derivative of the exponent at $t = 0$ gives the cumulant
\begin{equation}
  c_k = E_R\,\bigl\langle\omega^{k-1}\bigr\rangle\ (k\ \text{odd}) , \qquad c_k = E_R\,\Bigl\langle\omega^{k-1}\coth\frac{\omega}{2T}\Bigr\rangle\ (k\ \text{even}) ,
  \label{eq:D.32-ck}
\end{equation}
which at $T = 0$ is \cref{eq:8.5-cumulants}. The odd cumulants do not depend on temperature. For $k = 3$, \cref{eq:D.32-c3} agrees, since the on-site force constant of a harmonic crystal is $M_\kappa\langle\omega^2\rangle$ by completeness of the polarization vectors.
